% Root/B/D complete actual-file audits PASS; auxiliary, no object-atlas assertion.
\begin{proposition}[Proved: canonical energy on a minimal surface resolution]
\label{D:exceptional-canonical-energy}
Let $A$ be a normal Noetherian local $k$-algebra of dimension two,
with residue field the algebraically closed field $k$, and suppose that
$A$ admits a dualizing complex. Let
$\pi:Z\to\operatorname{Spec}A$ be a projective regular resolution. Suppose
that $p_g=\dim_kH^1(Z,\mathcal O_Z)<\infty$ and that no exceptional
curve is a smooth rational curve of self-intersection $-1$.
Write $E_1,\ldots,E_v$ for the exceptional integral curves, and set
\[
 b_i=-E_i^2,\quad g_i=p_a(E_i),\quad
 N=(-(E_i\cdot E_j))_{i,j},\quad
 \kappa_i=K_Z\cdot E_i=2g_i-2+b_i.
\]
When the exceptional set is empty, all the following sums and energies
are zero. Otherwise $N$ is positive definite with nonpositive
off-diagonal entries, $\kappa_i\ge0$, and
\[
 \mathcal E:=\kappa^{\mathsf t}N^{-1}\kappa
       \le 8p_g+2\sum_i b_i.                           \tag{D.CE.1}
\]
Let $\mathcal B$ be any union of exceptional components with
$\kappa_i=0$, and let $\mathcal C$ be its complement. Every member of
$\mathcal B$ is a smooth rational $(-2)$-curve. For nonempty
$\mathcal C$, form the Schur complement
\[
 S=N_{\mathcal C\mathcal C}
       -N_{\mathcal C\mathcal B}
        N_{\mathcal B\mathcal B}^{-1}N_{\mathcal B\mathcal C},
                                                               \tag{D.CE.2}
\]
with $S=N$ if $\mathcal B$ is empty. It is positive definite, has
nonpositive off-diagonal entries, and
\[
 \sum_{\substack{i\in\mathcal C\\g_i\ge1}}
       g_i^2(S^{-1})_{ii}
 \le \mathcal E,
 \qquad
 \sum_{\substack{i\in\mathcal C\\g_i\ge1}}
       \frac{g_i^2}{S_{ii}}
 \le \mathcal E.                                      \tag{D.CE.3}
\]
These are numerical assertions; they do not assume that the indicated
Schur complement has already been realized by a scheme contraction.
\end{proposition}
\begin{proof}
Negative definiteness is the usual exceptional intersection theorem.
Since an integral proper curve of arithmetic genus zero over the
algebraically closed residue field is a smooth rational curve,
minimality gives $b_i\ge2$ when $g_i=0$. Thus $\kappa_i\ge0$.
For $g_i\ge1$, we also have $\kappa_i\ge g_i$; this uses only
$b_i\ge1$ and does not assert the stronger inequality
$\kappa_i\ge2g_i$.

We recall explicitly the elementary matrix facts used below. For a
positive definite symmetric matrix with positive diagonal and
nonpositive off-diagonal entries, scale the diagonal to write it as
$D^{1/2}(1-P)D^{1/2}$ with $P$ symmetric and entrywise nonnegative.
Positive definiteness gives $\lambda_{\max}(P)<1$, and the largest
absolute eigenvalue of a nonnegative symmetric matrix is its largest
eigenvalue. Hence $(1-P)^{-1}=\sum_{m\ge0}P^m$ is entrywise
nonnegative. In particular $N^{-1}$ is entrywise nonnegative.
Put
\[
                      z_*=\frac{1}{2}N^{-1}\kappa\ge0.
\]
Choose a nearest nonnegative integer $z_i$ to each coordinate of
$z_*$, and write $\epsilon=z-z_*$, so $|\epsilon_i|\le1/2$.
For a nonzero effective exceptional cycle $D_z=\sum_i z_iE_i$,
the fibre-dimension bound gives
$H^2(Z,\mathcal O_Z(-D_z))=0$. The exact sequence for $D_z$
therefore gives a surjection
$H^1(Z,\mathcal O_Z)\to H^1(D_z,\mathcal O_{D_z})$ and hence
$p_a(D_z)\le p_g$. This is the same full, possibly nonreduced,
cycle argument as in Lemma~\ref{D:exceptional-weighted-genus}.
Cycle adjunction gives the exact identity
\[
 p_a(D_z)=1+\frac{\kappa^{\mathsf t}z-z^{\mathsf t}Nz}{2}
       =1+\frac{\mathcal E}{8}
                    -\frac{\epsilon^{\mathsf t}N\epsilon}{2}.
                                                               \tag{D.CE.4}
\]
Write $m_{ij}=E_i\cdot E_j$ for $i<j$. Positive definiteness on
the vector of all ones implies
$2\sum_{i<j}m_{ij}<\sum_i b_i$. Therefore
\[
 \epsilon^{\mathsf t}N\epsilon
 \le \frac{1}{4}\sum_i b_i+\frac{1}{2}\sum_{i<j}m_{ij}
 < \frac{1}{2}\sum_i b_i.                               \tag{D.CE.5}
\]
If $z\ne0$, (D.CE.4)--(D.CE.5) yield
$\mathcal E\le8p_g-8+2\sum_i b_i$, which implies (D.CE.1).

The rounded vector can be zero even when $\kappa\ne0$, so this
case must be handled separately. If $z=0$, then
$0\le(z_*)_i\le1/2$. Since the off-diagonal entries of $N$ are
nonpositive,
\[
 \frac{\mathcal E}{4}=z_*^{\mathsf t}Nz_*
       \le\sum_i b_i(z_*)_i^2
       \le\frac{1}{4}\sum_i b_i.
\]
This proves the stronger bound $\mathcal E\le\sum_i b_i$ in
this case, and finishes (D.CE.1).

If $\kappa_i=0$, adjunction and $b_i>0$ force $g_i=0,b_i=2$.
The Schur complement of a positive definite block is positive
definite. Its off-diagonal entries are nonpositive, because
$N_{\mathcal B\mathcal B}^{-1}$ is entrywise nonnegative and the
two rectangular off-diagonal blocks are entrywise nonpositive.
The block inverse identity gives
\[
 (N^{-1})_{\mathcal C\mathcal C}=S^{-1},\qquad
 \mathcal E=\kappa_{\mathcal C}^{\mathsf t}
                      S^{-1}\kappa_{\mathcal C}.
\]
The inverse $S^{-1}$ is entrywise nonnegative. Dropping its
off-diagonal contributions and using $\kappa_i\ge g_i$ proves
the first inequality of (D.CE.3). Finally
$(S^{-1})_{ii}\ge1/S_{ii}$: this follows from the scalar Schur
complement formula, or from Cauchy--Schwarz applied to
$S^{1/2}e_i$ and $S^{-1/2}e_i$. This proves the second inequality.
\end{proof}

\begin{corollary}[Proved: degree-linear effective genus energy on punctured covers]
\label{D:punctured-cover-linear-energy}
In the algebraic or descended strict-henselian punctured-cover setting
of Proposition~\ref{D:punctured-cover-linear-resolution}, replace the
chosen resolution by a minimal regular resolution of the covering
normal surface. There is a constant depending only on the fixed base
data such that, for every connected cover of degree $n$,
\[
 v+\sum_i b_i+\sum_i g_i+p_g+\mathcal E\le Cn.          \tag{D.CE.6}
\]
Every sum in (D.CE.3), for every choice of zero-canonical components
to eliminate, obeys the same kind of degree-linear bound.
\end{corollary}
\begin{proof}
The starting resolution has $v$, $\sum b_i$ and $\sum g_i$
bounded linearly in $n$ by
Proposition~\ref{D:punctured-cover-linear-resolution}; the canonical
defect theorem, Proposition~\ref{D:cover-uniform-canonical-defect},
gives the same bound for $p_g$, after completion if necessary.
Contract exceptional smooth rational $(-1)$-curves until none remains.
Existence and regularity of these contractions are the surface
contraction facts already used in constructing the preceding
resolution; see \citet[Lemma 54.16.8, Tag 0C2M]
{StacksDSurfaceResolution}. We do not require the resulting minimal
resolution to retain its morphism to the fixed base resolution.

For completeness, these contractions preserve all the needed
numerical bounds even when they destroy a normal-crossings
configuration. If the contracted curve meets the other exceptional
integral curves with intersection multiplicities $a_i$, their new
self-intersection and arithmetic genus are
\[
 (E_i')^2=E_i^2+a_i^2,
 \qquad p_a(E_i')=p_a(E_i)+\frac{a_i(a_i-1)}{2}.          \tag{D.CE.7}
\]
These follow from the total-transform formula
$\rho^*E_i'=E_i+a_iE$ and $K_Z=\rho^*K_{Z'}+E$.
The sum of the positive numbers $b_i$ drops by exactly
$1+\sum_i a_i^2$, while the sum of the arithmetic genera rises
by at most $\frac12\sum_i a_i^2$. Consequently the total genus
increase over the whole contraction sequence is at most half
the original $\sum b_i$. The number of components decreases,
and $p_g$ is unchanged under blowing down a regular $(-1)$-curve.
All four initial bounds therefore persist. Apply
Proposition~\ref{D:exceptional-canonical-energy} to finish.
\end{proof}

\begin{remark}[Scope of the energy estimate]
The diagonal numbers $S_{ii}$ in (D.CE.3) can be small rational
numbers after elimination of long rational blocks. The estimate
controls the corresponding genus contributions without a lower
bound on these individual numbers. It does not construct a
finite-type presentation of bundles or covers, control its
equation degrees, or prove an index-linear representation count.
The hereditary object-atlas problem remains separate.
\end{remark}
